\scanpage{148}
cette catégorie est semi-simple). Si alors $\mathcal C$ est une catégorie abélienne, les foncteurs additifs de $\mathcal M$ dans $\mathcal C$ sont autom. exacts. Cela prouve que le ``Pb universel''\inserted{ (avec grain de sel)} en termes de catégories abéliennes (Variante $2$) est résolu par $\mathcal M$.

Prouvons donc que $\mathcal M$ est une catégorie \uline{abélienne semi-simple}. Cela équivaut à ceci :

\deleted{tout morphisme est \unsure{sa}}

(i) Tout objet\deleted{\unsure{X}} satisfait condition des chaînes ascend\textsuperscript{tes}, et descendantes pour les sous-objets facteurs directs,

\deleted{i.e. l'ens des \unsure{fa}}

(ii) Tout morphisme $u:P\longrightarrow Q$ est ``direct'', i.e. isom à un morphisme
\[
 A\times B\xrightarrow[\operatorname{inj}_1\circ\operatorname{pr}_2]{}B\times C
\]

Pour vérifier (ii), \deleted{\unsure{supposons}} considérons $R=P\times Q$, et l'endomorphisme \deleted{$\unsure{\operatorname{pr}_2|\operatorname{pr}_1}$} $v=\operatorname{inj}_2\,u\operatorname{pr}_2$ de $R$. On voit aisément que pour \deleted{\unsure{que}} (ii)\inserted{ soit vrai} pour $u$, il\inserted{ \unsure{f.|faut} et} suffit \deleted{\unsure{de}} qu'il le soit pour $v$, donc on est ramené au cas où $P=Q$.
% slanted
\margin{donner lemme \unsure{commun|connexe} : il est innocent de modifier $P\to Q$ par
\diagram{\begin{tikzcd}[ampersand replacement=\&,column sep=small,row sep=small]
 P+P'\arrow[d] \& Q+Q'\arrow[dl,bend left=20]\\
 P\arrow[r] \& Q\arrow[u]
\end{tikzcd}}}
D'autre part, soit \deleted{\unsure{R}} $S$ un autre objet, et \unsure{sur|soit} $R=P+S$, $v=u\oplus\deleted{\unsure{1|\mathrm{id}}}0_S$, alors on constate que la propriété est vraie pour $u$ ssi elle l'est pour $v$. Cela nous
% archive sheet number: 147
